\documentclass[11pt,reqno]{amsart}
\usepackage{../notes/math327-lecture}
\newcommand{\rr}{\rightarrow}
\newtcolorbox{solutionbox}{lecturebox,colback=lightblue,colframe=courseblue,
  fonttitle=\bfseries,colbacktitle=lightblue,coltitle=courseblue,title=Solution}
\setlist[enumerate,1]{label=\arabic*.,leftmargin=*,itemsep=1em}
\setlist[enumerate,2]{label=(\alph*),leftmargin=*,itemsep=0.55em}
\pagestyle{plain}
\hypersetup{pdftitle={Midterm Exam Solutions, MATH 327, October 17, 2025},
  pdfauthor={Manish Patnaik},pdfsubject={MATH 327 past examination}}
\begin{document}
\begin{tcolorbox}[lecturebox,colback=lightblue,colframe=courseblue,
  boxrule=0pt,left=3mm,right=3mm,top=3mm,bottom=3mm]
  {\color{courseblue}\bfseries\Large Midterm Exam Solutions\hfill MATH 327}\par
  \smallskip
  \textbf{October 17, 2025}\hfill Instructor: Manish Patnaik
\end{tcolorbox}



\emph{You have 50 minutes for this exam. There is a total of \textbf{40 points} possible. If you referencing a result, please state the result you are using as clearly as possible. You may use any result proven in class except if it is essentially the same as what is being asked. You may not use results from another (algebra) class that we have not yet covered. If in doubt, please ask. Good luck! }  


\vspace{0.15in} 


\begin{enumerate}

\item (10 pts) Show that a group of order 21 must contain an element of order $3$

\begin{solutionbox}
	Let $G$ be such a group. It can have elements of order $1, 3, 7, 21$. If it has an element $x$ of order $21$, i.e., it is cyclic, then $x^7$ has order $3$. Indeed, $(x^7)^3=1$ and $(x^7)^2=x^{14} \neq 1.$ 
	So, next assume $G$ has no element of order $21$ and also, by way of contradiction, no elements of order $3$. So all the elements, other than the identity, have order $7$. Let $x, y$ be two such elements and $H_x$ and $H_y$  be the subgroups generated by $x$ and $y$. They each have seven elements. As $H_x \cap H_y$ is a subgroup of $H_x$ it has order either $1$ or $7$, i.e. $H_x$ and $H_y$ either are equal, or they meet only in the identity element. From these remarks, we see that $G$ is a union of subgroups $H_x$ for elements $x_1, x_2, \ldots, x_k$, and these subgroups overlap just at the identity. Hence $21 = 6k+1$. But this equation is impossible to solve for $k$ an integer.  
	
	
\end{solutionbox}
	
	
	%\item (5 pts) Let $G$ be a group and $H, K$ be proper subgroups (i.e $H, K \subsetneq G$). Show that $G$ cannot be the union of $H$ and $K$. 


\vspace{0.15in} 

\item (10 pts) Let $G$ be a finite group and consider the map $\varphi: G \rr G$ which sends $x \mapsto x^2$. Show that if $\varphi$ is an isomorphism if and only if $G$ is abelian and has odd order. (\textit{Note:} make sure to argue that $\varphi$ is a homomorphism as well.)

\begin{solutionbox}
For $\varphi$ to be a homomorphism, we must have $\varphi(ab)= \varphi(a) \varphi(b)$, i.e. $(ab)^2 = a^2 b^2 $. But from \be a (ba ) b = ab ab = a a b b = a( ab ) b \ee we conclude that $ab = ba$. So $\varphi$ is a homomorphism if and only if $G$ is abelian (as clearly if $G$ is abelian, this map is a homomorphism).

Now as $G$ is finite, the map $\varphi$ is an isomorphism if and only if it is injective, i.e. if and only if $x^2=1$ implies $x=1$. If $G$ has odd order, we cannot have any elements of order $2$ by Lagrange's theorem. Conversely, if $G$ has even order, then there must exist some element of order $2$: assume not, then we can write $G$ as a union of $1$ and subsets of size two $\{ y, y^{-1}\}$ forcing $G$ to have odd order. 

	
\end{solutionbox}

\vspace{0.15in} 


\clearpage
\item (10 pts) Let $H, K$ subgroups of a group $G$, and let $HK:= \{ hk \mid h \in H, k \in K \}.$ Assume that $H$ normalizes $K$, i.e. \be h K h^{-1} = K \mbox{ for all } h \in H. \ee
\begin{enumerate}
	\item (2pts) Show that $HK$ is a subgroup. 
	
	\begin{solutionbox}
		It suffices to show that $HK$ is closed under inverse and multiplication, as it clearly contains the identity. Let $h, h' \in H$ and $k, k' \in K$. Then \be h k h' k' = h h' (h'^{-1} k h') k' \ee and by the fact that $H$ normalizes $K$, the right hand side is again in $HK$. As for inverse, we note that \be (hk)^{-1} = k^{-1} h^{-1} \ee and also that $KH \subset HK$: indeed $k h = h (h^{-1} k h)$.
	\end{solutionbox}
	\item (2 pts) Show that $K$ is a normal subgroup of $HK$
	 and $H \cap K$ is a normal subgroup of $H$
	\begin{solutionbox} In the same notation as the previous part, we compute: 
		\be (hk) (k') (hk)^{-1} = hk k' k^{-1} h^{-1}  \in h K h^{-1} \subset K. \ee 
	Also, if $h \in H$ and $x \in H \cap K$, we have $h x h^{-1}$ lies in $K$ by the normalizing assumption, and also in $H$ since all three elements are in $H$.	
		
	\end{solutionbox}
	\item (6 pts) Prove that $H/H \cap K$ and $HK/K$ are isomorphic (construct a map, and argue that it is an isomorphism)
	
	\begin{solutionbox}
	Consider the map $\varphi:H \rr HK/K$ which sends $h \mapsto  h\cdot1K= hK$. It is a homomorphism since, by definition of the multiplication in $HK/K$, we have $(hK )(h' K) = h h' K.$ But the latter is just $\varphi(hh')$. 
	
	We contend that $\varphi$ is surjective. Indeed, given by $hkK$, we note that $\varphi(h) = hkK$, since $hkK= hK$.
	
	Finally, we argue that $\varphi$ has kernel $H \cap K$. Indeed, if $\varphi(h)=K$, then $hK=K$ which means that $h \in K$. By the first isomorphism theorem, we have $\varphi$ induces an isomorphism $H/H \cap K \rr HK/K$.		
	\end{solutionbox}
\end{enumerate}

\vspace{0.15in}

\clearpage
\item (10 pts) 

\begin{enumerate}
	\item (2 pts) For $G$ a group, show that conjugation defines a (left) group action of $G$ on itself, i.e. show that the following operation defines a left group action: \be \star: G \times G &\rr& G \\   g, x \mapsto g \star x &:=& g x g^{-1} \ee 

	\begin{solutionbox} We have that $ 1 \star x = 1 x 1^{-1} = x$. Also if $g, h \in G$, then \be (gh) \star x = (g h) x (gh)^{-1} = g h x h^{-1} g^{-1} = g (h x h^{-1}) g^{-1} = g \star ( h \star x) \ee \end{solutionbox}

	\item (2 pts) Show that the center of $G$, i.e. $Z(G):= \{ g \in G \mid gh =hg \mbox{ for all } h \in G \},$ is the union of all orbits of cardinality $1$. 
	
	\begin{solutionbox} If $z \in Z(G)$, then $g \star z = g z g^{-1} = z$ since $gz =zg$. So the orbit of $z$ has cardinality $1$, so $Z(G)$ is contained in the union of orbits of cardinality $1$. On the other hand, if $\{ x \}$ is an orbit (of cardinality $1$), this means that $g \star x =x$ for all $g \in G$, i.e. $gx = xg$ for all $g \in G$. Hence $x \in Z(G)$.
	
	\end{solutionbox}
	
	\item (2 pts) Show that a normal subgroup of $G$ must equal a union of orbits under this action. 
	
	\begin{solutionbox}
		If $N \subset G$ is normal and $n \in N$, then the orbits of $n$ must also be in $N$ as $g n g^{-1} \in N$ for all $n \in N$. 
		
		 	\end{solutionbox}
	
	\item(4 pts) Describe all the orbits and stabilizers for $G=S_3$ 
	
	\begin{solutionbox}
		There are three orbits, that of $(1)$, that of $(12)$ and that of $(123)$. The first orbit contains just one element, that of $(12)$ contains $\{ (12), (13), (23) \}$, and that of $(123)$ contains $\{ (123), (132) \}$. 
		
		The stabilizer of $(1)$ is all of $S_3$. The stabilizer of $(12)$ is $\{ 1, (12) \}$, and similarly for the other transpositions.
		
		The stabilizer of $(123)$ is $\{ 1, (123), (132) \}$ and this is the same for $(132)$.
	\end{solutionbox}
	
	
\end{enumerate}	





\end{enumerate}

\end{document}